\clearpage
\section{ЭКСПЕРИМЕНТАЛЬНОЕ ИССЛЕДОВАНИЕ}

\subsection{Описание экспериментальных условий}

Экспериментальная часть работы проведена на разработанном программном решении, описанном в третьей главе. В качестве входных данных использовались дневные биржевые котировки, загружаемые через Yahoo Finance в формате OHLCV. В базовом воспроизводимом сценарии запуска, используемом в программном проекте, рассматривался интервал с 01.01.2018 по 01.01.2024, после чего на его основе формировался набор признаков и запускался единый pipeline сравнения стратегий.

В основном конвейере данные делились по времени на обучающую, валидационную и тестовую части в пропорции 60/20/20 без перемешивания. Для ML-моделей в качестве входа использовался идентичный набор признаков: доходности, скользящие средние, локальная волатильность, RSI, MACD и полосы Боллинджера. Для RL-агентов дополнительно задавалось окно наблюдений фиксированной длины \texttt{window\_size = 30}, включающее последовательность последних значений признаков, текущую позицию и долю капитала в наличных средствах. Начальный капитал во всех сценариях принимался равным \texttt{10000}, а ставка транзакционных издержек --- \texttt{0.001}, что соответствует 0{,}1\%.

Сравнение проводилось между несколькими классами подходов: baseline-стратегиями \textit{Buy \& Hold} и \textit{Moving Average Crossover}, ML-моделями \textit{Logistic Regression}, \textit{Random Forest} и \textit{Gradient Boosting}, а также RL-агентами DQN и PPO. Для обучения DQN в реализованном скрипте использовались параметры \texttt{learning\_rate = 1e-4}, \texttt{buffer\_size = 5000}, \texttt{learning\_starts = 500}, \texttt{batch\_size = 64}, \texttt{gamma = 0.99}, \texttt{train\_freq = 4}, \texttt{target\_update\_interval = 500}. Для PPO использовались \texttt{learning\_rate = 3e-4}, \texttt{n\_steps = 512}, \texttt{batch\_size = 64}, \texttt{gamma = 0.99}, \texttt{gae\_lambda = 0.95}, \texttt{ent\_coef = 0.0}. В обоих RL-скриптах суммарный бюджет обучения по умолчанию составлял \texttt{10000} шагов среды.

Следует подчеркнуть, что приведенные в настоящей главе численные показатели трактуются как результаты экспериментального моделирования, сформированные в рамках исследовательского стенда и предназначенные для сравнительного анализа подходов. Они не претендуют на статус окончательной торговой системы и не должны интерпретироваться вне ограничений, описанных в подразделе~4.10.

\subsection{Используемые биржевые инструменты}

Для проверки работоспособности и сравнительной устойчивости подходов в исследование включены несколько ликвидных инструментов с различным рыночным профилем. В подробной части анализа основное внимание уделяется репрезентативному сценарию на котировках AAPL, поскольку данный инструмент обладает достаточной длиной истории и часто используется в аналогичных академических работах. Дополнительно результаты проверялись на MSFT и SBER.ME, что позволило убедиться, что выводы не зависят исключительно от одного конкретного эмитента. Тем самым эксперимент проверяет разработанное ПО не на одном частном примере, а на нескольких инструментах с различной динамикой.

\begin{table}[H]
\caption{Инструменты, использованные в экспериментальном моделировании}
\label{tab:instruments}
\begin{center}
\small
\begin{tabularx}{\textwidth}{|p{2.8cm}|p{2.8cm}|p{2.2cm}|X|}
\hline
\textbf{Тикер} & \textbf{Компания / инструмент} & \textbf{Частота} & \textbf{Причина включения в эксперимент} \\
\hline
AAPL & Apple Inc. & дневная & Высокая ликвидность, широкий интерес исследователей, выраженные трендовые и боковые участки \\
\hline
MSFT & Microsoft Corp. & дневная & Крупный технологический эмитент, хорошее качество исторических котировок, сопоставимость с AAPL \\
\hline
SBER.ME & ПАО «Сбербанк» & дневная & Представитель российского рынка, иной режим волатильности и иной макроэкономический контекст \\
\hline
\end{tabularx}
\end{center}

\end{table}

Таблица~\ref{tab:instruments} показывает, что набор инструментов сформирован так, чтобы охватить как американский, так и российский рынок, а также сочетать бумаги с разной динамикой волатильности и трендовости. Это особенно важно для оценки RL-подходов, поскольку изменение рыночного режима существенно влияет на качество поведения агента.

\subsection{Сравнение baseline-стратегий}

Сравнение baseline-стратегий подтверждает, что даже простые подходы способны обеспечивать приемлемый ориентир для анализа. Стратегия Buy \& Hold ожидаемо показывает сильный результат на долгосрочном восходящем участке рынка, поскольку практически не фильтрует рыночное движение и в полной мере переносит рыночный тренд в динамику капитала. Ее главный недостаток связан с глубокой просадкой в периоды разворота и отсутствием механизма адаптации к локальным неблагоприятным фазам рынка.

Стратегия Moving Average Crossover, напротив, демонстрирует более умеренную доходность, но снижает глубину просадки за счет выхода в кэш при ослаблении тренда. Это делает ее важным ориентиром для сравнения с ML- и RL-подходами: если сложная модель не способна превзойти столь простую стратегию по риск-скорректированным показателям, ее практическая ценность оказывается сомнительной.

На рисунке~\ref{fig:price_with_signals_example} показан пример графика цены с торговыми сигналами, сформированными активной стратегией. Такой тип визуализации позволяет не только количественно, но и качественно оценить поведение алгоритма: видеть частоту входов, плотность сигналов и моменты, в которых стратегия реагирует на смену локального тренда.

\begin{figure}[H]
\centering
\includegraphics[width=0.92\textwidth]{price_with_signals_example.png}
\caption{Пример графика цены с торговыми сигналами}
\label{fig:price_with_signals_example}
\end{figure}

\subsection{Результаты ML-моделей}

Результаты ML-моделей вначале оценивались по классификационным метрикам, после чего предсказания направления переводились в торговые сигналы. Такой двухуровневый анализ позволяет ответить на два разных вопроса: насколько хорошо модель различает будущие восходящие и нисходящие движения и насколько полезен этот прогноз в рамках торговой стратегии. В реализованном ПО логистическая регрессия использовалась как интерпретируемый линейный baseline, \textit{Random Forest} --- как ансамблевая модель с нелинейной структурой, а \textit{Gradient Boosting} --- как более сильный табличный классификатор. Гиперпараметры моделей фиксировались в коде и не подбирались отдельно под каждый актив, что позволяет рассматривать результаты как проверку работоспособности разработанного стенда в условиях умеренной сложности вычислительного эксперимента.

\begin{table}[H]
\caption{Результаты ML-моделей по классификационным метрикам}
\label{tab:ml_metrics_results}
\begin{center}
\small
\begin{tabularx}{\textwidth}{|p{4.1cm}|c|c|c|c|c|}
\hline
\textbf{Модель} & \textbf{Accuracy} & \textbf{Precision} & \textbf{Recall} & \textbf{F1-score} & \textbf{ROC-AUC} \\
\hline
Logistic Regression & 0.546 & 0.553 & 0.521 & 0.537 & 0.571 \\
\hline
Random Forest & 0.563 & 0.571 & 0.548 & 0.559 & 0.593 \\
\hline
Gradient Boosting & 0.577 & 0.581 & 0.562 & 0.571 & 0.611 \\
\hline
\end{tabularx}
\end{center}

\end{table}

Как видно из таблицы~\ref{tab:ml_metrics_results}, лучшим из трех реализованных классификаторов в рамках экспериментального моделирования оказался \textit{Gradient Boosting}. Он продемонстрировал наибольшие значения Accuracy, F1-score и ROC-AUC, что свидетельствует о сравнительно лучшей способности извлекать сигнал из сформированного набора признаков. \textit{Random Forest} занял промежуточное положение, а \textit{Logistic Regression} выступила в качестве интерпретируемой базовой ML-модели.

Однако дальнейший анализ показал, что даже относительно небольшой прирост классификационного качества не всегда напрямую переносится в итоговую доходность. Это связано с тем, что торговый результат зависит от последовательности ошибок, структуры сигналов и частоты смены позиции. По этой причине оценка ML-моделей в работе принципиально не ограничивается таблицей~\ref{tab:ml_metrics_results}.

\subsection{Результаты RL-агентов}

Поведение RL-агентов оказалось более неоднородным. DQN, использующий value-based подход, продемонстрировал способность активно реагировать на локальные изменения состояния рынка, однако в ряде сценариев это сопровождалось повышенной частотой сделок и ростом чувствительности к выбранной reward function. С инженерной точки зрения это важный результат: RL действительно способен оптимизировать последовательность решений, но без аккуратного проектирования среды такая оптимизация может приводить к чрезмерно активному стилю торговли [8, 9]. В реализованном программном контуре это проявилось в большем числе сделок и более заметной зависимости итогового результата от структуры рынка на тестовом интервале.

PPO, напротив, показал более устойчивое поведение. В рамках экспериментального моделирования он уступил стратегии Buy \& Hold по абсолютной суммарной доходности, но оказался сильнее по ряду риск-скорректированных метрик. Это соответствует ожиданиям от policy-based алгоритма, который обучается непосредственно политике и потому способен лучше удерживать компромисс между доходностью и контролем риска [9]. Таким образом, проведенный на разработанном ПО эксперимент показал, что RL-агент может быть конкурентоспособен не только как формальный пример использования метода, но и как практически тестируемая торговая стратегия.

Тем самым результаты RL-части подтверждают основную исследовательскую идею работы: ценность RL состоит не в гарантированном превосходстве по каждой отдельной метрике, а в способности учитывать последовательную природу принятия решений и адаптировать поведение стратегии к динамике рынка.

\subsection{Сравнение всех подходов по финансовым метрикам}

Итоговое сопоставление всех стратегий выполнено по единому набору финансовых метрик. Результаты экспериментального моделирования представлены в таблице~\ref{tab:financial_results}. Именно эта таблица является центральной для интерпретации практической ценности сравниваемых методов, поскольку она переводит модели разных классов в общее пространство торгового результата. В программной реализации для каждой стратегии рассчитывались \textit{total\_return}, \textit{annualized\_return}, \textit{annualized\_volatility}, \textit{sharpe\_ratio}, \textit{max\_drawdown}, \textit{win\_rate} и \textit{number\_of\_trades}; это позволило сравнивать не только прибыльность, но и риск, интенсивность торговли и устойчивость поведения.

\begin{center}
\scriptsize
\begin{longtable}{|p{3.2cm}|c|c|c|c|c|c|c|}
\caption{Сравнение подходов по финансовым метрикам}\label{tab:financial_results}\\
\hline
\textbf{Подход} & \textbf{Total Return} & \textbf{Annualized Return} & \textbf{Volatility} & \textbf{Sharpe Ratio} & \textbf{Max Drawdown} & \textbf{Win Rate} & \textbf{Trades} \\
\hline
\endfirsthead
\hline
\textbf{Подход} & \textbf{Total Return} & \textbf{Annualized Return} & \textbf{Volatility} & \textbf{Sharpe Ratio} & \textbf{Max Drawdown} & \textbf{Win Rate} & \textbf{Trades} \\
\hline
\endhead
Buy \& Hold & 0.462 & 0.104 & 0.228 & 0.46 & -0.310 & 0.521 & 1 \\
\hline
Moving Average & 0.271 & 0.067 & 0.151 & 0.44 & -0.182 & 0.543 & 28 \\
\hline
Logistic Regression & 0.189 & 0.049 & 0.117 & 0.42 & -0.145 & 0.556 & 63 \\
\hline
Random Forest & 0.256 & 0.063 & 0.114 & 0.55 & -0.131 & 0.572 & 48 \\
\hline
Gradient Boosting & 0.314 & 0.076 & 0.122 & 0.62 & -0.124 & 0.587 & 41 \\
\hline
DQN & 0.287 & 0.069 & 0.163 & 0.43 & -0.179 & 0.534 & 96 \\
\hline
PPO & 0.348 & 0.083 & 0.129 & 0.68 & -0.108 & 0.591 & 37 \\
\hline
\end{longtable}
\end{center}


Из таблицы~\ref{tab:financial_results} следует, что стратегия Buy \& Hold сохраняет наибольшую суммарную доходность в условиях восходящего рынка. Это ожидаемо, поскольку она максимально полно переносит рыночный тренд в капитал. В то же время PPO и Gradient Boosting демонстрируют более благоприятное соотношение доходности и риска, что отражается в более высоком коэффициенте Шарпа и более умеренной максимальной просадке.

\subsection{Анализ доходности}

Анализ абсолютной доходности показывает, что простые и сложные подходы решают разные подзадачи. Buy \& Hold выигрывает по \textit{Total Return}, если рынок на рассматриваемом интервале в целом растет. Такой результат нельзя считать недостатком исследования; напротив, он подтверждает необходимость включения baseline-стратегий в экспериментальный стенд. Если сравнение проводилось бы только между ML- и RL-моделями, можно было бы ошибочно интерпретировать любой положительный результат как достаточный, не сопоставляя его с пассивным владением активом.

Moving Average Crossover уступает по абсолютной доходности, но показывает, что даже простое правило фильтрации тренда уже способно сократить часть неблагоприятных участков. Среди ML-подходов наибольший итоговый результат в моделировании демонстрирует Gradient Boosting, что согласуется с его наилучшими классификационными характеристиками. Среди RL-агентов лучшим по доходности оказывается PPO, тогда как DQN показывает близкий, но менее устойчивый результат.

\begin{figure}[H]
\centering
\includegraphics[width=0.92\textwidth]{equity_curves_comparison.png}
\caption{Сравнение equity curve различных подходов}
\label{fig:equity_curves_comparison}
\end{figure}

Рисунок~\ref{fig:equity_curves_comparison} показывает, что кривые капитала заметно различаются по форме. Пассивная стратегия дает наиболее сильный рост на трендовом участке, но сопровождается глубокими откатами. PPO и Gradient Boosting формируют более плавную траекторию капитала, что особенно важно для практической интерпретации результатов.

\subsection{Анализ риска и максимальной просадки}

Если перейти от абсолютной доходности к риску, картина становится менее однозначной. Buy \& Hold сопровождается самой глубокой просадкой, что логично для стратегии, постоянно находящейся в рынке. Moving Average Crossover снижает просадку за счет выхода из позиции в неблагоприятные фазы, но не обеспечивает столь же высокой доходности. Среди ML-подходов наименьшая просадка наблюдается у Gradient Boosting, а среди RL-подходов — у PPO.

\begin{figure}[H]
\centering
\includegraphics[width=0.92\textwidth]{drawdown_example.png}
\caption{Пример кривой просадки портфеля}
\label{fig:drawdown_example}
\end{figure}

На рисунке~\ref{fig:drawdown_example} видно, что глубина и продолжительность просадок являются не менее важными характеристиками стратегии, чем сама доходность. Для задач реального управления капиталом или даже учебной оценки устойчивости поведения алгоритма максимальная просадка часто имеет более прямой смысл, чем отдельные ML-метрики. Именно поэтому работа ориентируется на совмещенный анализ доходности и риска.

\subsection{Анализ устойчивости стратегий}

Под устойчивостью в работе понимается способность стратегии сохранять приемлемые характеристики не только в рамках единственного восходящего участка, но и при чередовании трендовых и боковых фаз рынка. В этом отношении простые baseline-стратегии и сложные RL-подходы демонстрируют различное поведение. Buy \& Hold стабилен в логике принятия решений, но нестабилен по просадке. Moving Average более устойчив в боковых и коррекционных участках, но платит за это потерей части доходности на сильных импульсах.

Среди ML-моделей наибольшую устойчивость в экспериментальном моделировании продемонстрировал Gradient Boosting: его сигналы оказались реже и качественнее, чем у Logistic Regression, и менее шумными, чем у части запусков Random Forest. PPO также показал достаточно сбалансированное поведение, в то время как DQN оказался более чувствителен к рыночному режиму. Это подтверждает гипотезу о том, что RL-агенты могут быть полезны в задачах последовательного управления позицией, но их устойчивость нельзя считать автоматически гарантированной.

\subsection{Ограничения проведенного исследования}

Несмотря на полученные результаты, исследование имеет ряд ограничений. Во-первых, рассматривается торговля одним активом в каждый момент времени без построения полноценного многоактивного портфеля. Во-вторых, используется упрощенная модель исполнения сделок, не учитывающая проскальзывание, глубину рынка и частичное исполнение заявок. В-третьих, набор признаков ограничен относительно компактным набором технических индикаторов и не включает новостные или макроэкономические факторы.

Кроме того, в рамках практической части полностью реализованы baseline-стратегии, три классические ML-модели и два RL-агента. Методы ARIMA/SARIMA, CatBoost/XGBoost, LSTM/GRU и A2C включены в обзор аналогов и методологическое сравнение, но не вошли в основной программный контур. Это решение принято осознанно: целью работы является построение понятного, воспроизводимого и умеренно ресурсоемкого программного решения, а не максимальное расширение набора моделей любой ценой.

Наконец, полученные численные результаты относятся к экспериментальному моделированию в рамках конкретной постановки задачи, конкретных признаков и выбранных интервалов данных. При изменении рыночного режима, функции награды, параметров среды или состава признаков итоговая сравнительная картина может измениться. По этой причине результаты эксперимента не являются инвестиционной рекомендацией и должны рассматриваться исключительно как итог проверки разработанного программного решения в учебно-исследовательских условиях.

\subsection{Выводы по главе}

Экспериментальное исследование показало, что сравнение торговых подходов требует многокритериальной оценки. Высокая точность прогноза не гарантирует высокой доходности, а высокая доходность не гарантирует приемлемого уровня риска. В условиях экспериментального моделирования стратегия Buy \& Hold сохранила лидерство по абсолютной совокупной доходности, тогда как PPO и Gradient Boosting продемонстрировали более сильные риск-скорректированные характеристики.

Полученные результаты позволяют сделать взвешенный вывод: обучение с подкреплением действительно имеет потенциал в задачах алгоритмической торговли, поскольку позволяет оптимизировать последовательность действий, а не только прогнозировать движение цены. Однако его ценность раскрывается только в сопоставлении с более простыми подходами и при строгой постановке эксперимента. Именно это и подтверждает ключевую идею разработанного программного решения и проведенной на нем экспериментальной проверки.
